\section{Introduction}\label{sec:introduction}
% statements-of-record rows resolved here: none

What would a conditional closure of Strong Birch--Swinnerton-Dyer~\citep{BirchSwinnertonDyer1965,Tate1966bsd} look
like, and what exactly would it rest on?  The classical landscape already
contains powerful conditional and partial routes: Gross--Zagier and
Kolyvagin~\citep{GrossZagier1986,Kolyvagin1989} control the rank-one Heegner setting; Skinner--Urban and
Iwasawa-main-conjecture methods~\citep{SkinnerUrban2014} control large ordinary regions;
Cassels--Tate and Selmer-complex formalisms~\citep{Flach1990,Nekovar2006} control part of the
quadratic and determinant bookkeeping.  These results do not, by
themselves, give a single uniform proof of the strong formula in every
rank and at every bad prime.  The question addressed here is narrower
and more structural: if the remaining arithmetic inputs are isolated
with precise names, can one assemble them into a single audited closure
whose scalar readout is exactly the Strong-BSD identity?

The answer is deliberately conditional.  This is \emph{not} an
unconditional classical proof of BSD, scalar or strong.  The result is
theorem-grade within the Six Birds closure discipline, conditional on
the standing shell-closure record, the named recognition sources, and the
headline imported theorem stacks.  Outside the framework it reads as the
conditional theorem
\[
  \begin{aligned}
  &\SelBSD\text{-closure}
    \wedge \GammaPD \wedge \GammaSP \wedge \GammaGZ
    \wedge \chiCTp\text{-imports}
    \wedge \AofE\text{-imports} \\
  &\qquad {}\wedge\ \text{Beilinson-imports}
    \Longrightarrow \text{Strong BSD}.
  \end{aligned}
\]
This implication is the paper's result in its plainest form.  The work
below does not erase the hypotheses; it makes their roles explicit and
shows that, once they are supplied, the Strong-BSD scalar identity is
forced by a single closure architecture.

The first component of the architecture is a separation of the deep
arithmetic into two kinds of inputs.  The first kind consists of three
recognition sources.  A recognition source is named arithmetic content
carried by the formed BSD layer: it is supplied as part of the closure
object and used as a hypothesis, not derived from the general framework.
The three sources are \(\GammaPD\), the additive-prime \(p\)-adic
descent/Tamagawa correction; \(\GammaSP\), the signed-Selmer
persistence of the \(\Sha\)-factor through local cover change; and
\(\GammaGZ\), the higher-rank Gross--Zagier-type fixity of the
archimedean regulator readout.  They correspond, respectively, to the
local \(p\)-adic factor, the \(\Sha\)-persistence factor, and the
rank-\(\geq 2\) regulator/fixity factor in the Strong-BSD product.  In
the formal development they are typed carriers of hypotheses, not
free-standing assumptions of the framework; in the paper they are treated
as named recognition content and are introduced before they are used.

The second kind of input consists of imported theorem stacks.  The three
headline imports are the ones used directly in the landing theorem.  The
\(\chiCTp\) stack packages the Cassels--Tate Pfaffian comparison and its
sign convention.  The \(\AofE\) stack packages the Selmer-complex pairing
substrate needed to compare the determinant formulation with the
Cassels--Tate side.  The Beilinson stack packages the rank-\(\geq 2\)
archimedean determinant--\(L\)-value comparison.  Later sections also
record the overconvergent \(\eta\) and \(T_{\mathrm{CASCADE}}\) stacks as
supporting conditional imports, but those are not hidden inside the three
headline imports.  These imports name classical arithmetic geometry as
provenance: Sakamoto, Macias Castillo--Sano, Nekov\'a\v{r}, Flach, the
Cassels--Tate literature, Burns--Flach, Beilinson, and the surrounding
Selmer-complex and determinant technology.  They are not reproved here.
They enter as proof-carrying assumptions, and later sections keep their
individual roles separate.

The second component is the saturated BSD trace shell.  It is the
five-part object
\[
  \SelBSD=(E,\rhoE,J^{!}_{\mathrm{BSD}},
  \mathrm{Vis}^{!}_{\mathrm{BSD}},\mathrm{Audit}^{!}_{\mathrm{BSD}}),
\]
where \(E/\Q\) is the elliptic curve, \(J^{!}_{\mathrm{BSD}}\) records
the duality or involutive datum, \(\mathrm{Vis}^{!}_{\mathrm{BSD}}\)
records the visible scalar readout, and
\(\mathrm{Audit}^{!}_{\mathrm{BSD}}\) records the closure audit.  The
central scalar is the residual
\[
  \rhoE
  =
  \frac{L^{(r)}(E,1)}{r!}
  -
  \frac{\#\Sha(E/\Q)\,\Reg_{\mathrm{NT}}(E)\,\Omega_E
    \prod_v c_v(E)}
       {|E(\Q)_{\tors}|^2}.
\]
Thus \(\rhoE=0\) is not a new BSD-like statement: it is exactly the
standard Strong-BSD scalar identity written as a residual equation.
The point of the shell is to keep the analytic side, the arithmetic
product, the duality datum, the visible readout, and the audit data in
one typed object, so that later arguments can say precisely which input
changes which part of the scalar package.

The recognition predicate \(\PiBSD\) is the shell's structured readout.
It is not defined to be Strong BSD.  It is a quantified predicate that
asks for the three recognition-source readouts in their proper ranges:
the \(p\)-adic descent readout for every additive bad prime and
trivializing local lift in scope, the signed-Selmer persistence readout
for every admissible signed-Selmer prime, and the higher-Gross--Zagier
fixity readout in rank at least two, with the classical low-rank lane
for ranks \(0\) and \(1\).  The readout-level equivalence says that
vanishing of the shell residual is equivalent to the usual Strong-BSD
scalar identity.  This equivalence is a genuine algebraic step: it
rewrites the residual, rather than making \(\PiBSD\) a synonym for the
conjecture.

The headline landing then has a short conceptual chain.  The recognition
sources and imports are supplied.  The higher-rank fixity source forces
the shell residual to vanish, after the \(\kapr\)-normalization identifies
the cascade regulator product with the standard
\(\#\Sha/|E(\Q)_{\tors}|^2\) factor.  The readout equivalence converts
\(\rhoE=0\) into
\[
  \frac{L^{(r)}(E,1)}{r!}
  =
  \frac{\#\Sha(E/\Q)\,\Reg_{\mathrm{NT}}(E)\,\Omega_E
    \prod_v c_v(E)}
       {|E(\Q)_{\tors}|^2}.
\]
The master-theorem applicability audit checks that the three recognition
sources have the correct framework profile, and the composite signature
checks that the three source-governed factors are not disposable: changing
the \(\Sha\)-factor, the regulator factor, or the Tamagawa factor changes
the scalar package.  This is the sense in which the conclusion is a
closure result rather than a restatement of its target.

Two ingredients are imported from \citet{TsiokosBSDApparatus2026}.
First, the \(\kapr\) normalization in that work is the bridge from the cascade leading-term
form to the standard Strong-BSD factor \(\#\Sha/|E(\Q)_{\tors}|^2\).
Second, its adequacy and five-column framing provide the
diagnostic language for separating height, period, finite, \(p\)-adic,
and determinant contributions.  This article uses those results as cited
inputs; it does not refer to internal labels of the cited work.  The AOR
terminology used at the end of the paper is the Audited Operational
Realisability language of \cite{TsiokosAOR2026}, and the recognition-mode
discipline is part of the audited interaction calculus of
\cite{TsiokosFoundationsIII2026}.

The final move is a reading of the completed carrier as a refinement-stable
audited object.  Audited Operational Realisability is a refinement-stable,
audit-closed membership criterion: informally, an object qualifies when
the residuals left by refinement have been accounted for by approved
mechanical records, imported records, or bridged recognition records.  In
this paper the object is \(\SelBSD\).  The resulting statement is the
framework's ``Strong BSD holds in reality'' reading, but it is
non-eliminative: the recognition records are reclassified as bridged
records, not made to disappear.

We close the introduction by fixing the scope.  This is not an
unconditional classical proof of BSD.  It is a conditional closure under
the standing shell-closure record, the three recognition sources, and the
headline imported theorem stacks displayed above.  The rank-\(\geq 2\)
generalized Gross--Zagier identity
\(T_{E12}\) and the additive \(p\leq 3\) Iwasawa main conjecture
\(T_{\mathrm{BAD}}\) remain open audit-result no-gos: the paper records
where they are needed and how the recognition sources carry them, but
does not prove them.  Likewise, the AOR-instance reading is
non-eliminative; it records an audit-closed presentation of the conditional
closure, not the discharge of every open arithmetic obligation.
